
\begin{thebibliography}{*****}




\bibitem[AL1]{AL-horseshoe} A. Avila and M. Lyubich. 
  The full renormalization horseshoe for unimodal maps of higher
  degree: exponential contraction along hybrid classes.
Publ. Math. IHES 114 (2011), 171-223. 

\bibitem[AL2]{AL-posmeas} A. Avila and M. Lyubich. 
Lebesgue measure of Feigenbaum Julia sets. arXiv:1504.02986. 


\bibitem[BC]{BC} X. Buff and A. Cheritat. Quadratic Julia sets of
  positive area. Annals Math., v. 176 (2012), 673--746.

\bibitem[BD]{BD} B. Branner and A. Douady. Surgery on complex polynomials. Holomorphic dynamics, {M}exico, 1986.


\bibitem[BR]{BR}L. Bers and H. L. Royden. Holomorphic families of injections. Acta Math. 157, 259--286 (1986).


 \bibitem[Che]{Che:nearpar} A. Cheritat. Near parabolic renormalization for unicritical holomorphic maps. arXiv:1404.4735.


   
 \bibitem[Chi]{Ch} D. K.~Childers. Are there critical points on the boundaries of mother hedgehogs? Holomorphic dynamics and renormalization, Fields Inst. Commun., 53 (2008), 75–87.  



\bibitem[D1]{D-Siegel} A. Douady. Disques de Siegel et anneaux de Herman. In S\'eminaire Bourbaki,
  Ast\'erisque, v. 152--153 (1987), 151--172. 



\bibitem[D2]{D-discont} A. Douady. Does a Julia set depend continuously on a polynomial?
   Proceed. Symp. in Applied Math., v. 49 (1994), 91--139.

\bibitem[DH1]{DH:Orsay} A. Douady and J. H. Hubbard.
  \'{E}tude dynamique des polyn\^{o}mes complexes.
 Publication Mathematiques d'Orsay,
84-02 and 85-04.

\bibitem[DH2]{DH}
 A. Douady and J. H. Hubbard. On the dynamics of polynomial-like maps.
   Ann. Sc. \'{E}c. Norm. Sup., v. 18 (1985), 287 -- 343.


\bibitem[DL]{DL} D. Dudko and M. Lyubich. Local connectivity of the Mandelbrot set at some satellite parameters of bounded type. arXiv:1808.10425.


\bibitem[GJ]{GJ} J. Graczyk and P. Jones. Dimension of the boundary of quasiconformal Siegel disks. Invent. math. 148, 465--493 (2002).


\bibitem[GY]{GY} D. Gaidashev and M. Yampolsky. Renormalization of almost commuting pairs. arXiv:1604.00719.


\bibitem[H]{H} M. Herman. Conjugaison quasi symm\'etrique des diff\'eomorphisms du cercle \`a des rotations
et applications aux disques singuliers de Siegel. Manuscript (1986).

\bibitem[HPS]{HPS} M. W. Hirsch, C. C. Pugh, M. Shub. Invariant Manifolds. Lecture Notes in Mathematics, v. 583. Springer-Verlag, Berlin-New York, 1977.






\bibitem[IS]{IS} H. Inou and M. Shishikura.
  The renormalization for parabolic fixed points and their
  perturbations. Manuscript 2008. 

\bibitem[La]{La}  P. Lavaurs. Syst\`emes dynamiques holomorphes:
explosion des points p\'eriodiques paraboliques. The\`se, Univirsit\'e
Paris-Sud, 1989. 


\bibitem[L1]{L:FCT} M. Lyubich. Feigenbaum-Coullet-Tresser Universality and Milnor's Hairiness Conjecture.
   Ann. Math., v. 149 (1999), 319 - 420.

\bibitem[L2]{L-book} M. Lyubich. Conformal Geometry and Dynamics of
  Quadratic Polynomials. Book in preparation, www.math.sunysb.edu/~mlyubich/book.pdf. 



\bibitem[McM1]{McM2} C. McMullen. Renormalization and three manifolds
  which fiber over the circle.   Princeton University Press, 1996. 

\bibitem[McM2]{McM3} C. McMullen. Self-similarity of Siegel disks and Hausdorff dimension of Julia sets.
  Acta Math., v. 180 (1998), 247--292.  



\bibitem[MN]{MN} N. S. Manton and M. Nauenberg. Universal scaling behaviour for iterated maps in the complex plane.
Commun. Math. Phys. 89(1983), 555-570.


\bibitem[MP]{MP} R. S. MacKay and I. C. Persival. Universal small-scale structure near the boundary of Siegel disks of
arbitrary rotation numer. Physica 26D(1987), 193-202.



\bibitem[Pe]{Pe} C. Petersen. Local connectivity of some Julia sets containing a circle with an irrational rotation, Acta. Math.,177, 1996, 163-224.


\bibitem[PM]{PM}  R. P\'erez-Marco. Fixed points and circle maps, Acta Math., 179 (1997), 243-294.


\bibitem[R]{R} J. Riedl. Arcs in Multibrot Sets, Locally Connected Julia Sets and Their Construction by Quasiconformal Surgery. Ph.D. thesis, TU M\"unchen, 2000. IMS Thesis Server,  http://www.math.stonybrook.edu/ims-thesis-server

\bibitem[S]{S} D. Sullivan. Bounds, quadratic differentials and renormalization conjetures.
  AMS Centennial Publ. 2: Mathematics into Twenty-first century, 1992. 


\bibitem[St]{St} A. Stirnemann. Existence of the Siegel disc renormalization fixed point. Nonlinearity 7(1994), no. 3, 959--974.

\bibitem[ST]{ST} D. Sullivan and W. Thurston. Extending holomorphic motions. Acta Math. 157, 243--257 (1986).


\bibitem[Sw]{Sw} G. Swiatek. On critical circle homeomorphisms. Bol. Soc. Bras. Mat., v. 29 (1998), 329--351.





\bibitem[Wi]{Wi} M. Widom. Renormalization group analysis of quasiperiodicity in analytic maps. Commun. Math. Phys. v. 92 (1983),121--136.

 




\bibitem[Ya]{Ya-posmeas} M. Yampolsky.  Siegel Disks and Renormalization Fixed Points.
  Fields Inst. Comm., v. 53 (2008), 377--393.  




\end{thebibliography}
